\section{Stress Tests That Catch the Degeneracy}
\label{sec:stress}

% STRETCH SECTION. Target: 1.25 pages (~1000 words).
% Decision gate: end of week 2. If §5 experiments aren't landed by then, CUT
% this section, expand §5 to fill the page budget, and move this material to
% appendix as "Future work / proposed evaluations."

\todo{This section is optional (stretch \S6 per the experiment checklist). Go/no-go
by end of week 2 based on whether \S\ref{sec:blindness} experiments are landed.
If cut, move this content to the appendix.}

\subsection{Pattern-distribution divergence}
\label{sec:stress:pdiv}

% Source: E7. Run SISC on the synthetic output, get the (pattern, length,
% magnitude) histogram, compare to the real-data histogram via JSD or KL.
% Threshold: the real-data distribution is roughly uniform over K patterns;
% any synthetic generator concentrated on <3 patterns fails.

\subsection{Pattern-transition entropy}
\label{sec:stress:tentropy}

% Source: E9. Compute H(p_{t+1} | p_t) over the synthetic pattern chain.
% On real data: typically > 1 bit. On a collapsed PEM: near 0 bits.
% This is the diagnostic that most directly catches the PEM argmax
% collapse described in §4.

\subsection{Volatility autocorrelation decay}
\label{sec:stress:vol}

% Source: E10. Fit an exponential decay to the autocorrelation of |returns|.
% The paper claims this should be slowly decaying ("stylized fact"); a linear
% ramp has near-zero autocorrelation at all lags. Quantify and table.

\subsection{Validation: the stress tests on a second generator}
\label{sec:stress:validation}

% Source: E11. Apply the three metrics to TimeGAN's output. Argue they
% distinguish "degenerate" from "non-degenerate" generators without being
% specific to FTS-Diffusion's particular failure mode.

\begin{table}[t]
  \centering
  \todo{Table: proposed metrics evaluated on real data, FTS-Diffusion (ours),
  TimeGAN, linear-ramp. Columns: metric, real, FTS-Diff, TimeGAN, ramp.}
  \caption{Proposed stress-test metrics across generators. \todo{Fill from
  E9--E11.}}
  \label{tab:stress}
\end{table}
